% Part: first-order-logic
% Chapter: natural-deduction
% Section: provability-consistency

\documentclass[../../../include/open-logic-section]{subfiles}

\begin{document}

\iftag{FOL}
      {\olfileid{fol}{ntd}{prv}}
      {\olfileid{pl}{ntd}{prv}}

\olsection{\usetoken{S}{derivability} lan Konsistensi}

Saiki kita bakal netepake sawetara sipat relasi !!{derivability}.
Sipat kasebut narik kawigaten dhewe-dhewe, nanging saben sipat bakal
duwe peran ing bukti teorema kelengkapan.

\begin{prop}\ollabel{prop:provability-contr}
  Yen $\Gamma \Proves !A$ lan $\Gamma \cup \{!A\}$ ora konsisten,
  mula $\Gamma$ ora konsisten.
\end{prop}

\begin{proof}
Ayo !!{derivation} saka~$!A$ adhedhasar~$\Gamma$ dijenengi~$\delta_1$,
lan !!{derivation} saka~$\lfalse$ adhedhasar $\Gamma \cup \{!A\}$
dijenengi~$\delta_2$. Banjur kita bisa !!{derive}:
\begin{prooftree}
\AxiomC{$\Gamma, \Discharge{!A}{1}$}
\RightLabel{$\delta_2$}
\DeduceC{$\lfalse$}
\DischargeRule{\Intro{\lnot}}{1}
\UnaryInfC{$\lnot !A$}
\AxiomC{$\Gamma$}
\RightLabel{$\delta_1$}
\DeduceC{$!A$}
\RightLabel{\Elim{\lnot}}
\BinaryInfC{$\lfalse$}
\end{prooftree}
Ing !!{derivation} anyar iki, asumsi~$!A$ wis !!{discharged}, mula iki
!!a{derivation} adhedhasar~$\Gamma$.
\end{proof}

\begin{prop}
\ollabel{prop:prov-incons}
$\Gamma \Proves !A$ yen lan mung yen $\Gamma \cup \{\lnot !A\}$ ora
konsisten.
\end{prop}

\begin{proof}
Sepisan, upamana $\Gamma \Proves !A$, yaiku ana
!!a{derivation}~$\delta_0$ saka~$!A$ adhedhasar asumsi
!!{undischarged}~$\Gamma$. Kita entuk !!a{derivation} saka $\lfalse$
adhedhasar $\Gamma \cup \{\lnot !A\}$ kaya mangkene:
\begin{prooftree}
  \AxiomC{$\lnot !A$}
  \AxiomC{$\Gamma$}
  \RightLabel{$\delta_0$}
  \DeduceC{$!A$}
  \RightLabel{\Elim{\lnot}}
  \BinaryInfC{$\lfalse$}
\end{prooftree}

Saiki upamana $\Gamma \cup \{\lnot !A\}$ ora konsisten, lan ayo
$\delta_1$ dadi !!{derivation} kang cocog saka~$\lfalse$ adhedhasar
asumsi !!{undischarged} ing~$\Gamma \cup \{\lnot !A\}$. Kita entuk
!!a{derivation} saka~$!A$ adhedhasar $\Gamma$ wae kanthi
nggunakake~$\FalseCl$:
\begin{prooftree}
  \AxiomC{$\Gamma, \Discharge{\lnot !A}{1}$}
  \RightLabel{$\delta_1$}
  \DeduceC{$\lfalse$}
  \RightLabel{\FalseCl}
  \DischargeRule{\FalseCl}{1}
  \UnaryInfC{$!A$}
\end{prooftree}
\end{proof}

\begin{prob}
Buktekna yen $\Gamma \Proves \lnot !A$ yen lan mung yen $\Gamma \cup
\{!A\}$ ora konsisten.
\end{prob}

\begin{prop}\ollabel{prop:explicit-inc}
  Yen $\Gamma \Proves !A$ lan $\lnot !A \in \Gamma$, mula $\Gamma$
  ora konsisten.
\end{prop}

\begin{proof}
  Upamana $\Gamma \Proves !A$ lan $\lnot !A \in \Gamma$. Banjur ana
  !!a{derivation}~$\delta$ saka~$!A$ adhedhasar~$\Gamma$. Gatekna
  aplikasi prasaja aturan $\Elim{\lnot}$ iki:
  \begin{prooftree}
    \AxiomC{$\lnot !A$}
    \AxiomC{$\Gamma$}
    \RightLabel{$\delta$}
    \DeduceC{$!A$}
    \RightLabel{\Elim{\lnot}}
    \BinaryInfC{$\lfalse$}
  \end{prooftree}
  Amarga $\lnot !A \in \Gamma$, kabeh asumsi !!{undischarged} kalebu
  ing~$\Gamma$, mula iki nuduhake yen $\Gamma \Proves \lfalse$.
\end{proof}

\begin{prop}\ollabel{prop:provability-exhaustive}
  Yen $\Gamma \cup \{!A\}$ lan $\Gamma \cup \{\lnot !A\}$ kalorone
  ora konsisten, mula $\Gamma$ ora konsisten.
\end{prop}

\begin{proof}
Ana !!{derivation}s $\delta_1$ lan $\delta_2$ saka~$\lfalse$ adhedhasar
  $\Gamma \cup \{ !A \}$ lan saka $\lfalse$ adhedhasar $\Gamma \cup
  \{ \lnot !A \}$, miturut urutane. Banjur kita bisa !!{derive}
\begin{prooftree}
\AxiomC{$\Gamma, \Discharge{\lnot !A}{2}$}
\RightLabel{$\delta_2$}
\DeduceC{$\lfalse$}
\DischargeRule{\Intro{\lnot}}{2}
\UnaryInfC{$\lnot \lnot !A$}
\AxiomC{$\Gamma, \Discharge{!A}{1}$}
\RightLabel{$\delta_1$}
\DeduceC{$\lfalse$}
\DischargeRule{\Intro{\lnot}}{1}
\UnaryInfC{$\lnot !A$}
\RightLabel{\Elim{\lnot}}
\BinaryInfC{$\lfalse$}
\end{prooftree}
Amarga asumsi $!A$ lan $\lnot !A$ wis !!{discharged}, iki
!!a{derivation} saka~$\lfalse$ adhedhasar~$\Gamma$ wae. Mula $\Gamma$
ora konsisten.
\end{proof}

\end{document}
